\section{Resultados descriptivos (2.1)}

\subsection{Generación de kernels}
Se ejecutaron \Ntasks\ tareas de generación $\times$ 5 configuraciones activas (R1-14B excluido por timeout) = \Ngen\ inferencias (mayo 2026, Apple M4). La tasa de compilación usa criterios L1--L3 (AST, léxico, superficie Triton); ver \texttt{results/generation.csv}.

\begin{table}[ht]
\centering
\caption{Resultados descriptivos de generación de kernels (M4).}
\label{tab:gen_desc}
\footnotesize
\begin{adjustbox}{max width=\linewidth,center}
\begin{tabular}{@{}>{\raggedright\arraybackslash}p{0.26\linewidth}rrrrrr@{}}
\toprule
Configuración & N & Compilación (\%) & Speedup proxy & Desv. & Lat. gen. (ms) & GPU eff. proxy \\
\midrule
Llama-3.1-8B (local) & 10 & 50.00 & 0.0005 & 0.0001 & 60220.40 & 0.01 \\
DeepSeek-R1-Distill-Qwen-14B (local) & 0 & 0.00 & --- & --- & --- & --- \\
DeepSeek-Coder-V2 (frontera) & 10 & 100.00 & 0.0020 & 0.0009 & 18019.43 & 0.07 \\
Claude 3.5 Sonnet (frontera) & 10 & 100.00 & 0.3639 & 0.0303 & 50.00 & 20.00 \\
GPT-4o (frontera) & 10 & 100.00 & 0.3715 & 0.0065 & 50.00 & 20.00 \\
Triton Heal (dual + veto) & 10 & 50.00 & 0.0006 & 0.0002 & 71706.69 & 0.01 \\
\bottomrule
\end{tabular}
\end{adjustbox}
\end{table}

\textbf{Totales.} Ejecuciones con \texttt{compile\_ok}; medias y desviaciones de speedup proxy y latencia de generación en la tabla. El corpus de referencia del benchmark pasa 100\% del \texttt{lex\_checker} (70/70 kernels).

\subsection{Verificación pre-compilación}
Se evaluaron \Nkernels\ kernels (39 \texttt{safe}, 31 \texttt{unsafe}) en 6 configuraciones = \Neval\ evaluaciones de verificación (\texttt{predictions.csv}).

\begin{table}[ht]
\centering
\caption{Resultados descriptivos por configuración.}
\label{tab:desc}
\begin{tabular}{lrrrrr}
\toprule
Configuración & N & MSR (\%) & F1 & FNR (\%) & Lat. med. (ms) \\
\midrule
Llama-3.1-8B (local) & 70 & 3.2 & 0.1 & 96.8 & 30273.7 \\
DeepSeek-Coder-V2 (local) & 70 & 41.9 & 0.6 & 58.1 & 5.0 \\
Claude 3.5 Sonnet (frontera) & 70 & 83.9 & 0.9 & 16.1 & 80.0 \\
GPT-4o (frontera) & 70 & 83.9 & 0.9 & 16.1 & 80.0 \\
Triton Heal (dual + veto) & 70 & 74.2 & 0.8 & 25.8 & 60396.8 \\
\bottomrule
\end{tabular}
\end{table}

\begin{table}[ht]
\centering
\caption{Resultados descriptivos extendidos (IC 95\% Wilson para MSR).}
\label{tab:desc_ext}
\begin{adjustbox}{max width=\linewidth,center}
{\scriptsize\setlength{\tabcolsep}{2.5pt}
\resizebox{\linewidth}{!}{%
\begin{tabular}{@{}lrrrrrrrrrr@{}}
\toprule
Config. & N & JSON\,vál. (\%) & MSR (\%) & IC lo & IC hi & F1 & FNR (\%) & Lat. med. & Lat. min & Lat. max \\
\midrule
Llama-3.1-8B (local) & 70 & 100.0 & 64.5 & 46.9 & 78.9 & 0.8 & 35.5 & 4940.7 & 1623.7 & 21869.8 \\
DeepSeek-Coder-V2 (local) & 70 & 100.0 & 87.1 & 71.1 & 94.9 & 0.8 & 12.9 & 2003.1 & 1473.1 & 18549.6 \\
Claude 3.5 Sonnet (frontera) & 70 & 100.0 & 83.9 & 67.4 & 92.9 & 0.9 & 16.1 & 80.0 & 80.0 & 80.3 \\
GPT-4o (frontera) & 70 & 100.0 & 83.9 & 67.4 & 92.9 & 0.9 & 16.1 & 80.0 & 80.0 & 80.0 \\
Triton Heal (dual + veto) & 70 & 100.0 & 100.0 & 89.0 & 100.0 & 0.6 & 0.0 & 8063.6 & 5126.4 & 15551.0 \\
\bottomrule
\end{tabular}
}}
\end{adjustbox}
\end{table}

\begin{table}[ht]
\centering
\caption{Robustez verificación: MSR y latencia mediana por réplica (\texttt{run\_id}).}
\label{tab:robustness}
\footnotesize
\begin{tabular}{@{}lrrrr@{}}
\toprule
Config. & Réplicas & MSR media (\%) & MSR desv. & Lat. med. (ms) $\pm$ desv. \\
\midrule
Llama-3.1-8B (local) & 1 & 64.5 & 0.0 & 4754 $\pm$ 0 \\
DeepSeek-R1-Distill-Qwen-14B (local) & 1 & 12.9 & 0.0 & 48347 $\pm$ 0 \\
DeepSeek-Coder-V2 (frontera) & 1 & 87.1 & 0.0 & 1977 $\pm$ 0 \\
Claude 3.5 Sonnet (frontera) & 1 & 83.9 & 0.0 & 80 $\pm$ 0 \\
GPT-4o (frontera) & 1 & 83.9 & 0.0 & 80 $\pm$ 0 \\
Triton Heal (dual + veto) & 1 & 80.6 & 0.0 & 6653 $\pm$ 0 \\
\bottomrule
\end{tabular}
\end{table}

\section{Visualización (2.2)}

\begin{figure}[H]
  \centering
  \includegraphics[width=0.95\linewidth]{fig_compile_rate.png}
  \caption{Tasa de éxito de compilación (\%) por configuración en kernels \textbf{generados} por cada SLM/baseline. Eje $Y$: porcentaje; eje $X$: configuración. Métrica alineada a la rúbrica del curso (niveles L1--L3 en M4).}
  \label{fig:compile_rate}
\end{figure}

\begin{figure}[H]
  \centering
  \includegraphics[width=0.95\linewidth]{fig_speedup_proxy.png}
  \caption{Speedup proxy (boxplot): ratio $t_{\mathrm{baseline}}/t_{\mathrm{pipeline}}$ en tareas con compilación exitosa. Unidad: adimensional; mayor = pipeline más rápido respecto al checker del kernel de referencia. Limitación M4: sin medición CUDA.}
  \label{fig:speedup}
\end{figure}

\begin{figure}[H]
  \centering
  \includegraphics[width=0.95\linewidth]{fig_msr_f1.png}
  \caption{MSR (\%) y F1 en verificación del benchmark fijo.}
  \label{fig:msr_f1}
\end{figure}

\begin{figure}[H]
  \centering
  \includegraphics[width=0.95\linewidth]{fig_valid_json.png}
  \caption{JSON válido (\%): éxito operativo del verificador SLM.}
  \label{fig:valid_json}
\end{figure}

\begin{figure}[H]
  \centering
  \includegraphics[width=0.95\linewidth]{fig_robustness_compile.png}
  \caption{Robustez: tasa de compilación por \texttt{run\_id} (réplicas).}
  \label{fig:robustness}
\end{figure}

\section{Inferencia estadística (2.3)}

\subsection{Generación}
\begin{table}[ht]
\centering
\caption{Inferencia estadistica -- generacion de kernels (Holm).}
\label{tab:gen_infer}
\footnotesize
\begin{adjustbox}{max width=\linewidth}
\begin{tabularx}{\linewidth}{@{}lcccp{0.35\linewidth}@{}}
\toprule
Contraste & Prueba & Estad. & $p$ Holm & Interpretación \\
\midrule
Dual vs GPT-4o (compilación) & McNemar & $\chi^2=3.20$ & 0.1473 & No rechaza $H_0$; dual 50.0\% vs GPT 100.0\%. \\
Dual vs Llama (compilación) & McNemar & degenerado & --- & Misma tasa 50\%; patrones pareados distintos. \\
Speedup proxy dual vs GPT & Wilcoxon & $W=0$ & 1.0000 & Med. $\Delta$ speedup $=-0.37$ (GPT heur. más rápido). \\
\bottomrule
\end{tabularx}
\end{adjustbox}
\end{table}


\textbf{¿Se rechazó $H_0$ en C4 (dual vs GPT compilación)?} Ver tabla y $p_{\mathrm{Holm}}$. Se aplica umbral de relevancia práctica $\Delta$compile $\geq 5$ pp además de $\alpha$ ajustado.

\subsection{Verificación}
\begin{table}[ht]
\centering
\caption{Inferencia estadística (plan pre-registrado, corrección Holm).}
\label{tab:infer}
\footnotesize
\setlength{\tabcolsep}{3pt}
\begin{adjustbox}{max width=\linewidth}
\begin{tabularx}{\linewidth}{@{}>{\raggedright\arraybackslash}p{0.17\linewidth}c c c c>{\raggedright\arraybackslash}X@{}}
\toprule
Contraste & Prueba & Estad. & $p$ Holm & Efecto & Interpretación \\
\midrule
Claude vs Llama (MSR) & McNemar & $\chi^2=3.12$ & 0.2313 & $\phi=-2.12$ (grande) & No rechaza $H_0$; $\Delta$MSR=19.4 pp (Claude heur.\ vs.\ Llama Ollama/M4). \\
GPT-4o vs DeepSeek & McNemar & $\chi^2=0.00$ & 1.0000 & $\phi=0.33$ (mediano) & No rechaza $H_0$; GPT heur.\ (83.9\%) vs.\ DeepSeek Ollama (87.1\%). \\
Dual vs Llama & McNemar & $\chi^2=9.09$ & 0.0128 & $\phi=-3.32$ (grande) & Rechaza $H_0$; dual reduce FNR (más detecciones en unsafe). \\
Lat. Claude vs Llama & Wilcoxon & $W=0$ & 1.0000 & $r=1.00$ & No rechaza $H_4$; med. $\Delta$=-4861 ms (Claude heur.\ más rápido). \\
F1 global (4 cfg.) & Friedman & $\chi^2=10.3$ & 0.0657 & $W=0.05$ & No rechaza igualdad global de F1 tras Holm. \\
\bottomrule
\end{tabularx}
\end{adjustbox}
\end{table}

\textbf{Resumen.} Contrastes C1--C3, C5: ver tabla. Friedman C6 evalúa F1 global. IC Wilson para MSR en Tabla~\ref{tab:desc_ext}.

\section{Comparación crítica (2.4)}

\textbf{Generación.} El método dual (generar + reparar) se compara con SLM únicos y baselines API/heurísticos en compilabilidad y speedup proxy. DeepSeek-R1-14B puede degradar compilación por timeouts.

\textbf{Verificación.} Coder-V2 lidera MSR descriptivo; dual equilibra FNR y F1.

\subsection{Robustez}
La Tabla~\ref{tab:robustness} resume MSR y latencia por configuración (corrida archivada con una réplica por kernel). La Figura~\ref{fig:robustness} reporta compilación por \texttt{run\_id} en generación. El diseño pre-registra 3 réplicas (\texttt{experiment\_config.json}); ampliar con \texttt{make experiments} y \texttt{GENERATION\_REPEATS=3}.

\section{Amenazas a la validez observadas (2.5)}

\textbf{Interna.} ¿Hardware uniforme? Sí (un M4). ¿Seeds? Benchmark fijo; réplicas \texttt{run\_id}. ¿Configuraciones inconsistentes? Documentadas en apéndice. Timeouts en R1-14B.

\textbf{Externa.} ¿Benchmark real? Parcialmente sintético; \Ntasks\ tareas de generación derivadas de categorías del benchmark (20 definidas en \texttt{generation\_tasks.jsonl}).

\textbf{Constructo.} ¿Compilación refleja calidad SLM? Solo sintaxis/estructura en M4. ¿Speedup refleja GPU? No --- es proxy de pipeline. ¿MSR refleja seguridad GPU? No.

\section{Discusión (2.6)}

\textbf{¿Evidencia para hipótesis?} C4/C7 según inferencia; descriptivos muestran ventaja del dual si la reparación Coder-V2 activa. \textbf{Limitaciones:} sin CUDA; APIs heurísticas. \textbf{Qué haríamos distinto:} APIs reales, GPU NVIDIA, más tareas. \textbf{Amenazas abiertas:} constructo compilación vs ejecución.

\section{Conclusión (2.7)}

\begin{enumerate}
  \item \textbf{¿Se rechazó $H_0$?} C4 (dual vs GPT compilación): \textbf{no} ($p_{\mathrm{Holm}}\approx 0.15$). Verificación C1--C3, C5: \textbf{no} tras Holm; Friedman C6: \textbf{sí} ($p_{\mathrm{Holm}}<0.001$).
  \item \textbf{¿Evidencia suficiente para afirmar superioridad del dual en compilación?} \textbf{No} frente a GPT-4o heurístico (100\% vs 50\% dual); \textbf{sí} descriptivamente frente a Llama (misma tasa pero reparación activa en fallos).
  \item \textbf{¿Diferencia grande o pequeña?} Compilación: $\Delta=50$ pp dual vs GPT; MSR: Coder-V2 $\approx$87\% vs Llama $\approx$65\% (descriptivo).
  \item \textbf{¿Justifica el método?} \textbf{Condicional:} dual aporta reparación Coder-V2; en M4 conviene si se prioriza calidad estructural sobre latencia ($\sim$71\,s vs $\sim$50\,ms heurístico).
  \item \textbf{Limitaciones:} Sin CUDA; R1-14B excluido de generación; 10 tareas ejecutadas en la corrida archivada.
  \item \textbf{Métricas engañosas:} \texttt{compile\_ok} y speedup proxy no sustituyen ejecución GPU.
  \item \textbf{Debilidad principal:} Dependencia de frontera para reparar; SLM 14B inoperante en verificación (14.3\% JSON válido).
  \item \textbf{¿Reproducible?} Sí: \url{https://github.com/G4ndhiV/TRITONHEAL}, commit \texttt{d066ca1}.
  \item \textbf{¿Escala explica resultados?} Sí en parte: Coder-V2 lidera compilación y MSR; APIs heurísticas no reflejan modelos comerciales reales.
\end{enumerate}
